\section{Solutions for Chapter~\ref{ch:overlaps}}

\paragraph{Exercise~\ref{ex:overlap-phases}.}
\phantomsection
\label{sol:ex:overlap-phases}
The phase of $\lambda^3$ is
$3\theta_\phi+3\theta_\chi-3\theta_\psi$. Multiplication by
$c_\psi$ and division by $c_\phi c_\chi$ cancel it. The same
bookkeeping for $\nu^3$ gives
$3\theta_\phi+3\theta_\psi+3\theta_\eta$, exactly the
denominator phase. For $\kappa^3c_\eta$ it gives
$3\theta_\chi+3\theta_\psi-3\theta_\eta+3\theta_\eta$,
again exactly the denominator phase. The definitions also give
\[
 \frac{I_\nu I_\kappa}{I_\lambda}
 =\frac{\nu^3\kappa^3}{\lambda^3c_\psi^3}=J^3.
\]
All these ratios require their displayed denominators to be nonzero.
If a symmetry maps each field to a unit phase times its adjoint,
it conjugates each three-point coefficient. All unit phases
cancel in $I_\lambda$, so invariance says
$I_\lambda=\overline{I_\lambda}$. No claim that individual
coefficients are real was required.



\paragraph{Exercise~\ref{ex:overlap-descendant}.}
\phantomsection
\label{sol:ex:overlap-descendant}
The commutator identity
\[
 L_1L_{-1}^n|h\rangle
 =n(2h+n-1)L_{-1}^{n-1}|h\rangle
\]
follows by commuting $L_1$ past the $n$ lowering operators.
The sum of their $2L_0$ terms is
$2\sum_{j=0}^{n-1}(h+j)=n(2h+n-1)$.
Writing $N_n=\|L_{-1}^n|h\rangle\|^2$ therefore gives
$N_n=n(2h+n-1)N_{n-1}$, with $N_0=1$, and hence
$N_n=n!(2h)_n$. The first three
terms of the scalar vacuum-source series are
\[
 1+\Delta^2+
 \left[\frac{\Delta(\Delta+1)}2\right]^2+\cdots.
\]
The equal-time integral near coincidence contains
$\int_0^\varepsilon\theta^{-2\Delta}\,d\theta$, finite
precisely for $2\Delta<1$. At the endpoint it diverges
logarithmically.



\paragraph{Exercise~\ref{ex:overlap-cumulant}.}
\phantomsection
\label{sol:ex:overlap-cumulant}
Treat the normalized overlaps as a linear moment functional
with $\langle1\rangle=1$, and set $\widetilde X_i=X_i-m_i$.
Then
\[
 B_{ij}=\langle\widetilde X_i\widetilde X_j\rangle,
 \qquad K=\langle\widetilde X_1\widetilde X_2\widetilde X_3\rangle.
\]
After $X_i\mapsto u_iX_i+v_i$, the new centered variable is
$u_i\widetilde X_i$. Multilinearity proves the factors
$u_iu_j$ and $u_1u_2u_3$, including for complex data. The
third leg's $u_3$ denotes its amplitude factor after converting
the outgoing ket change into a bra change.

A physical branch insertion does not invalidate these affine
identities. It changes the moment functional itself, replacing
three-point moments by normalized four-point correlators.
Their connected part contains extra OPE data. Thus it invalidates
the inference that the invariant equals the original three-field
coefficient, unless an additional Ward or selection argument
removes that branch contribution.



\paragraph{Exercise~\ref{ex:overlap-filter}.}
\phantomsection
\label{sol:ex:overlap-filter}
Each retained-state summand is nonnegative. For every omitted
state, its exponential factor is at most
$e^{-2\tau(E_{\rm cut}-E_\phi)}$. Factoring this out of the
omitted sum proves the bound. To see why residuals do not provide
the cutoff, take a diagonal Hamiltonian with energies $0,1,2$
and retain the exact eigenvectors at $0$ and $2$. Both residuals
vanish, but the omitted state has energy $1<2$. A source with
weight on all three detects the failure of using the largest
retained energy as an omitted-state threshold.



\paragraph{Exercise~\ref{ex:overlap-response}.}
\phantomsection
\label{sol:ex:overlap-response}
For $H(\theta)=H+\theta V$, Hellmann--Feynman gives
$d_a=\partial_\theta g_a=\langle a|V|a\rangle-\langle0|V|0\rangle$
along the chosen differentiable eigenstate branch. The quotient rule gives
\[
 \partial_\theta\!\left(\frac{g_a}{g_b}\right)
 =\frac{d_ag_b-g_ad_b}{g_b^2}.
\]
In a degenerate space one must first choose the branches by diagonalizing
the perturbation there. Under $V\mapsto V+a1+bH$, the gap derivatives
become $d_a+bg_a$, so the added numerator is
$bg_ag_b-bg_ag_b=0$. A coupling reparametrization multiplies
every coordinate derivative by the same Jacobian. Comparing
its direction with size drift assumes the perturbation probes
the same running continuum coupling and that other corrections
are negligible. The algebraic cancellation of energy rescaling
does not prove either premise.



\paragraph{Exercise~\ref{ex:overlap-current}.}
\phantomsection
\label{sol:ex:overlap-current}
Write $u'_i=\sum_k u_kU_{ki}$ and
$v'_j=\sum_l v_lV_{lj}$. The gluing amplitude is linear in
both incoming kets, giving $M'=U^TMV$. Both $U^T$ and $V$
are unitary, so $M'^\dagger M'$ is unitarily conjugate to
$M^\dagger M$ and has the same eigenvalues.

For normalized pairs, the tensor-product difference satisfies
\[
 \|u\otimes v-u'\otimes v'\|
 \leq\|(u-u')\otimes v\|+
      \|u'\otimes(v-v')\|\leq\delta_1+\delta_2.
\]
Taking an inner product with a unit parent cannot increase that
norm. Divide by $|A_{000}|$, subtract the continuum gluing error
$\eta$, and use $|x^TMy|\leq\|M\|_2$ for normalized coefficient
vectors. This proves the stated bound. A small denominator
amplifies an absolute vector error; cancellation of a parent's
overall normalization is not cancellation of its directional
state error.



\paragraph{Exercise~\ref{ex:overlap-branch}.}
\phantomsection
\label{sol:ex:overlap-branch}
Expanding either side gives
$2b^4-4b^3+6b^2-4b+2$. Consequently
\[
 \sum_{i<j}\frac{\delta B_{ij}}{B_{ij}}
 =\frac{2N_\Lambda}{c}
   \frac{1+(b-1)^4+b^4}{b^4(b-1)^4}
 =2\frac{\delta K}{K}.
\]
This is exactly the cancellation in the logarithmic variation
of the estimator. Multiplying the identity in $q,r$ by
$b^4(b-1)^4$ reduces it to the same polynomial identity.
The quadratic variation is a bilinear Hessian $H(v_b,v_b)$;
differentiation controls $H(v_b,\partial_bv_b)$. It supplies
no identity for $H(v_b,v_{b'})$ at independent $b,b'$, so two
separately regulated singularities require another calculation.


